\documentclass[11pt]{article}
\usepackage[T1]{fontenc}
\usepackage[utf8]{inputenc}
\usepackage{lmodern,microtype}
\usepackage[margin=1in,headheight=15pt]{geometry}
\usepackage{amsmath,amssymb,amsthm,mathtools}
\usepackage{booktabs,array,longtable,enumitem}
\usepackage{xcolor,hyperref,fancyhdr,listings}
\definecolor{inkblue}{RGB}{27,61,87}
\hypersetup{colorlinks=true,linkcolor=inkblue,citecolor=inkblue,urlcolor=inkblue,
 pdftitle={Connected Causal Kernels and Commuting-Layer Search for Reidemeister-III Unlocking},
 pdfauthor={Research contribution prepared for the ProveIt project}}
\pagestyle{fancy}\fancyhf{}
\fancyhead[L]{\small Connected causal kernels for RIII unlocking}
\fancyhead[R]{\small ProveIt research note}
\fancyfoot[C]{\thepage}
\setlength{\parskip}{0.35em}
\setlength{\parindent}{1.25em}
\setlength{\emergencystretch}{2em}
\setlist{itemsep=0.2em,topsep=0.4em}
\lstset{basicstyle=\ttfamily\footnotesize,breaklines=true,columns=fullflexible,
 frame=single,framerule=0.3pt,xleftmargin=0.4em,xrightmargin=0.4em,
 aboveskip=0.8em,belowskip=0.8em,showstringspaces=false}
\newtheorem{theorem}{Theorem}[section]
\newtheorem{lemma}[theorem]{Lemma}
\newtheorem{proposition}[theorem]{Proposition}
\newtheorem{corollary}[theorem]{Corollary}
\theoremstyle{definition}
\newtheorem{definition}[theorem]{Definition}
\newtheorem{remark}[theorem]{Remark}
\newcommand{\poly}{\operatorname{poly}}
\newcommand{\RI}{\mathrm{RI}}
\newcommand{\RII}{\mathrm{RII}}
\newcommand{\RIII}{\mathrm{RIII}}
\newcommand{\supp}{\operatorname{supp}}
\newcommand{\code}[1]{\texttt{#1}}
\newcommand{\repo}{\url{https://github.com/VladimirReshetnikov/ProveIt}}
\begin{document}
\begin{titlepage}
\vspace*{0.15in}
{\large\sffamily\color{inkblue} TOPOLOGY / UNKNOT RECOGNITION\par}
\vspace{0.25in}
{\LARGE\bfseries Connected Causal Kernels and\par
Commuting-Layer Search for\par
Reidemeister-III Unlocking\par}
\vspace{0.3in}
{\large Single-exponential local search and a sharper route\par
toward quasi-polynomial unknot recognition\par}
\vspace{0.25in}
{\large Research contribution prepared for Vladimir Reshetnikov's ProveIt project\par}
\vspace{0.15in}
{8 October 2026\par}
\vspace{0.2in}
\begin{abstract}
We study the exact local subproblem of deciding whether at most $k$
Reidemeister-III moves can expose a crossing-decreasing Reidemeister-I or
Reidemeister-II move in an $N$-crossing classical knot diagram. The inspected
ProveIt implementation bounds a related support-aware search by
$N^{b+O(1)}(Ck)^k$, where $b$ is the number of independent births.
We first replace chronological interleavings by commuting layers and obtain
$N^{b+O(1)}2^{O(k)}$. A state-dependent version of the usual heap construction
admits an injective encoding by bounded-slot forests; the exact forest count is
$\frac{b}{j}\binom{dj}{j-b}$.

A second argument removes the birth parameter for the unlocking problem.
The backwards causal cone of one terminal reduction has a set of at most
$3k+2$ core crossings connected in the \emph{initial} projection graph,
even though the intervening moves change that graph. Full read/write
footprints are retained. Enumerating these connected cores and applying
the layer search gives a deterministic
$\poly(N,k)2^{O(k)}$ algorithm, with polynomial workspace, for every bounded
unlocking instance. A deliberately loose explicit bound is
$\poly(N,k)2^{23k+11}$.

This strengthens a component of the recognition pipeline, not the general
recognition theorem. If the actual deterministic simplification run has
unlocking budget $g$ and residual core size $r$, exact recognition costs
$\poly(N,g)2^{O(g)}+2^{O(r)}$; thus $g,r=O(\log^2 N)$ suffice for the
specific target $N^{O(\log N)}$. No uniform such bounds are established.
The accompanying standard-library Python kernels, certificate verifier,
27 unit tests, exhaustive Boolean comparisons, and actual-diagram audits
make the local claims reproducible. Paired timings show fewer interleavings
but no demonstrated wall-clock improvement on the measured knot kernels;
no production default change is recommended.
\end{abstract}
\vfill
{\small\sffamily Status: self-contained mathematical proofs and tested research code.
Not a formally verified or externally peer-reviewed result. Global priority
of the new local-search bounds has not been established.\par}
\end{titlepage}
\begingroup
\setlength{\parskip}{0pt}\small
\tableofcontents
\endgroup
\newpage

\section{Scope, baseline, and the contribution}
\label{sec:scope}

The computational target is exact recognition of a knot from a classical
one-component diagram, with an eventual worst-case bound of
$N^{O(\log N)}$. Here $N$ denotes the number of crossings; logarithms in
complexity comparisons are to base two. In the broader literature,
``quasi-polynomial'' can mean $\exp((\log N)^{O(1)})$. We retain the more
specific $\exp(O(\log^2 N))$ target used in this project.

The relevant inspected repository material is the integrated causal RIII
search and its audit in \code{synthesis/causal\_r3.tex}, together with
\code{fast/fastunknot/causal\_r3.py}, the dart operations in
\code{simplify.py}, and the diagram conventions in
\code{diagram.py}~\cite{repo-causal,repo-code,repo-simplify,repo-diagram}.
The implementation reads and analysis are recorded in \code{SOURCES.json}.
The directly pinned source revision is
\begin{center}\small\ttfamily
34f2689e96e04a13a32d133173b5e41a34320086.
\end{center}
The source files fetched just before this pin are additionally identified by
blob hashes; they are not silently treated as a complete checkout.

The existing causal method distinguishes an initial collection of independent
birth moves from later moves intersecting accumulated support. Its analysis
has the shape $N^{b+O(1)}(Ck)^k$. It also explains why physical-state-only
memoization and indiscriminate inverse pruning can destroy completeness with
a birth constraint. Its integrated trace representation identifies an actual
triangular face rather than merely a triple of crossings~\cite{repo-causal}.
Those are sound prerequisites to preserve, not replace.

Our first change is to enumerate \emph{layers of commuting moves}, removing a
factorial-scale interleaving overhead. Our second change uses the fact that
unlocking has a \emph{local terminal witness}: a specific empty monogon or
eligible bigon. This allows us to search a connected support of size linear
in $k$, rather than paying $N^b$ for unrelated starting regions.

\begin{center}
\begin{tabular}{>{\raggedright\arraybackslash}p{0.27\textwidth}>{\raggedright\arraybackslash}p{0.32\textwidth}>{\raggedright\arraybackslash}p{0.31\textwidth}}
\toprule
Claim & Scope & Status here\\
\midrule
Commuting-layer bound & General bounded-local action systems with the stated
commutation contracts & Proved; executable enumeration\\
Connected-support bound & Classical knot diagrams; RIII moves followed by one
RI/RII test & Proved; executable region search\\
Quasi-polynomial recognition & Actual run has $g,r=O(\log^2N)$ & Conditional
corollary, not a uniform theorem\\
Production speedup & Current complete \code{fastunknot} pipeline & Not measured;
no default switch proposed\\
\bottomrule
\end{tabular}
\end{center}

The main theorem is stronger than a promise about a small birth number:
\emph{every} instance of the bounded unlocking problem is covered. It is
weaker than general recognition because a diagram can require a much larger
unlocking budget, or the permitted crossing-nonincreasing strategy can stall.
These quantifiers must not be conflated.

\subsection{Relation to established work}

Heaps of pieces and Cartier--Foata normal forms are classical ways to identify
words differing only by commutations~\cite{viennot,krattenthaler}. The
contribution here is not the invention of those normal forms. It is the
explicit treatment of state-dependent local actions, the bounded-slot count,
the dart constants, and the connected initial-support argument needed for
this recognition pipeline.

Legrand-Duchesne, Rai, and Tancer study the number of moves required to
\emph{untangle the whole diagram}, including the parameter called defect,
and establish W[P]-completeness for that parameter~\cite{lrt}. Our parameter
counts only RIII moves in one episode ending in \emph{one} reduction. Their
parameter, their terminal condition, and their permitted global histories
are different. In particular, local fixed-parameter tractability proved here
does not decide whether an arbitrary diagram admits a globally optimal
untangling sequence.

Lackenby's July 2026 preprint develops an incompressibility and hierarchy
algorithm with an application to the unknot~\cite{lackenby}. Section 9
discusses the remaining cost analysis and the precision of its algorithmic
steps. The present paper does not implement those hierarchy operations or
supply their missing estimates. It addresses a rigorously specified local
kernel in the existing alternative pipeline. For a complete fallback we may
use reduced Khovanov homology, whose unknot detection is a theorem of
Kronheimer and Mrowka~\cite{km}; Bar-Natan's computational work provides a
principal antecedent for practical complex simplification~\cite{bn}.

\section{Bounded unlocking and the local action model}
\label{sec:model}

\begin{definition}[Bounded RIII unlocking]
Given a connected classical one-component diagram $D$ and a nonnegative
integer $k$, the problem $\mathrm{Unlock}_{\RIII}(D,k)$ asks whether a legal
sequence of at most $k$ RIII moves produces a diagram admitting a
crossing-decreasing RI or RII move. The terminal reduction need not yet be
performed. There are no crossing-increasing moves or crossing deletions
within the RIII sequence.
\end{definition}

An initially available reduction gives a zero-move witness. An empty diagram
has no crossing-decreasing move and therefore has a negative answer to this
\emph{local} question; the outer recognizer of course knows that the empty
knot diagram is an unknot. A completed negative local answer says only that
the prescribed episode does not exist. A resource cap gives an inconclusive
answer even to the local question.

For the abstract layer theorem, use a fixed finite site set $X$. A state
assigns local data to those sites, possibly subject to global invariants. A
legal action $a$ at a state has a unique key, a nonempty conservative footprint
$F(a)\subseteq X$, and a transition to a new state. The key is a stable
identifier of that particular local move, not merely a name for the set of
vertices it involves.

\begin{definition}[Locality and commutation contracts]
The action system satisfies the following requirements.
\begin{enumerate}[label=(\roman*)]
\item The footprint contains all sites whose data determine legality, all
data read to perform the action, and all data changed by the action.
\item Disjoint actions preserve one another's legality, keys, footprints,
and effects. They commute as transitions. In a legal chronological trace,
adjacent occurrences with disjoint footprints can be swapped in either
direction, retaining the same occurrence data.
\item Legal actions and their footprints are deterministically enumerable.
State descriptions, enumeration, key comparison, transitions, and target
tests when used have costs polynomial in the input size and in the
trace-length bound.
\item A bound $d$ is available such that, after a layer of disjoint actions,
at most $d$ legal actions intersect the footprint of any specified action
of that preceding layer.
\end{enumerate}
\end{definition}

Requirement (ii) is deliberately stronger than saying that two initially
legal operations happen to have the same final result in one example. It
includes the exchange property for occurrences in a trace. Read/write
locality with stable footprint descriptions supplies this property. An
implementation using an insufficient read set does not satisfy the theorem.

One general way to obtain (iv) is to bound every footprint by $f$ sites and
every site's incidence with legal action footprints by $B$. Then $d=fB$
is valid. We will obtain a better RIII-specific value $d=36$.

\begin{definition}[Births and dependency heights]
For a chronological trace $a_1\cdots a_m$, write $F_i=F(a_i)$ for the
footprint at its occurrence. Occurrence $j$ is a birth when
$F_j\cap\bigcup_{i<j}F_i=\varnothing$. Its dependency height is
\[
 h_j=1+\max\bigl(\{h_i:i<j,\ F_i\cap F_j\ne\varnothing\}\cup\{0\}\bigr).
\]
The $\ell$th layer is the set of occurrences having height $\ell$, sorted
by their action keys.
\end{definition}

Notice the distinction between accumulated support and the previous layer.
The old birth-front search branches on the former. The new normal form
requires every action after the first layer to meet the \emph{immediately
preceding} layer. That stronger condition is a consequence of dependency
height, not a heuristic restriction to the neighborhood of the last move.

\section{State-dependent commuting layers}
\label{sec:layers}

\begin{theorem}[Dynamic layer normal form]
\label{thm:foata}
Under the locality and commutation contracts, every legal trace has a legal
layered representative with the same length, endpoint, and occurrence
footprints. Actions within each layer are pairwise footprint-disjoint.
Every action in layer $\ell>1$ meets an action in layer $\ell-1$. The first
layer consists exactly of the original births. Conversely, a legal sequence
of pairwise-disjoint layers with this previous-layer condition is the
height-layer representative of its own chronological expansion.
\end{theorem}

\begin{proof}
If $i<j$ and $F_i\cap F_j\ne\varnothing$, then $h_j\ge h_i+1$.
Consequently two occurrences of equal height have disjoint footprints.
If $h_j=\ell>1$, the maximum in the defining recurrence is attained at an
earlier occurrence of height $\ell-1$. Thus the claimed previous-layer
intersection exists.

Orient every intersecting pair of occurrences from its earlier member to
its later member. This directed graph is acyclic because chronological
indices increase along its edges. The original word is a linear extension
of its transitive closure. Ordering occurrences first by height and then by
key is another linear extension. Two linear extensions of a finite poset
can be connected by adjacent swaps of incomparable elements: move the first
element of the desired order leftwards across its predecessors in the
current order, and repeat on the remaining suffix. Incomparable occurrences
here have disjoint footprints. Requirement (ii) therefore makes each swap
legal and preserves every involved footprint and key, including when the
action catalogue is state-dependent. The sorted word is a legal trace with
the same endpoint.

Height one means that there is no earlier intersecting occurrence, exactly
the definition of birth. Since equal-height keys are distinct in their
common pre-layer state, the canonical order is unambiguous.

For the converse, proceed by layers. Actions in the first layer have no
earlier intersecting action. An action in the next layer has an intersecting
predecessor in the previous layer and no earlier action at a greater
height. Equal-layer actions are disjoint. Its height is consequently exactly
its layer index. Induction proves the assertion.
\end{proof}

\subsection{What the normal form does and does not identify}

The normalization removes permutations of independent occurrences. It does
not identify all words with the same physical endpoint: inverses, different
local relations, and different histories can still lead to the same state.
The algorithm is not a shortest-path quotient by physical-state equality.
This is intentional, since support history can matter to the restricted
birth problem.

For example, take four Boolean sites initially zero and the following legal
involutions, as in the repository's inverse-pruning audit~\cite{repo-causal}:
\begin{center}
\begin{tabular}{llll}
\toprule
Action & Flipped bits & Guard & Footprint\\
\midrule
$A$ & $0,1$ & $x_2=x_3=0$ & $\{0,1,2,3\}$\\
$B$ & $2$ & $x_0=0$ & $\{0,2\}$\\
$C$ & $3$ & $x_1=0$ & $\{1,3\}$\\
\bottomrule
\end{tabular}
\end{center}
For the terminal predicate $x_2=x_3=1$, the trace $BC$ has two births,
whereas $AABC$ has one. Its layers are $\{A\},\{A\},\{B,C\}$.
Deleting the neutral pair $AA$ changes the birth count. The implementation
therefore retains immediate inverses. The root-free local algorithm also
retains them, avoiding a second set of pruning assumptions.

\subsection{Enumeration}

At the initial state enumerate pairwise-disjoint nonempty subsets of legal
actions, optionally imposing a first-layer size bound $b$. At subsequent
states, enumerate the legal actions meeting the previous layer's footprint
union, then enumerate their nonempty pairwise-disjoint subsets. Charge a
layer its number of actions, not one unit. Stop extending at total charge
$k$ and test the target at every completed prefix, including the empty one.

Testing only at completed layers is complete for bounded reachability.
Starting with any successful chronological prefix, Theorem~\ref{thm:foata}
produces a completed-layer trace reaching the same endpoint. There is no
need to insert arbitrary partial executions of a chosen layer into that
trace's search path: those smaller layers are already independently
enumerated by the algorithm.

An explicit depth-first stack holds one suspended subset iterator per
active level. It never stores a global table of all reached states.
Resource exhaustion propagates as an exception distinct from complete
bounded failure. The reference implementation uses immutable states, so a
failed or interrupted search cannot corrupt its input diagram.

\section{The bounded-slot forest count}
\label{sec:forests}

The crude chronological estimate gives $O(k)$ locally available extensions
at each of $k$ stages and hence a factor $(Ck)^k$. Merely replacing that
estimate by a verbal appeal to commutation would be insufficient. We give
an explicit encoding that also handles state-dependent actions.

\begin{lemma}[Forest encoding]
\label{lem:encoding}
Fix an initial layer of $b\ge1$ actions. Layered traces with $j$ occurrences
in total inject into ordered forests with $b$ rooted trees, where every
node has $d$ distinguished child slots, each either empty or occupied by
one child. The same upper bound covers attempted last layers rejected for
internal footprint overlap, provided all preceding layers are legal.
\end{lemma}

\begin{proof}
After a legal layer, order its actions by key. Every candidate for the next
layer intersects at least one parent action. Assign it to the first such
parent. For each parent, sort its assigned candidates by key and number
them from one. There are at most $d$ of them by contract (iv). For each
selected child, record the corresponding slot at its assigned parent.
Repeat for all layers; the roots are the prescribed initial actions in
canonical order. Each nonroot has one parent in the previous layer, so the
result is an ordered $d$-slot forest.

To prove injectivity, reconstruct breadth-first. The root layer and input
state determine the first post-layer state. That state determines the
candidate catalogue, parent assignments, and slot meanings. The occupied
slots uniquely determine the next selected actions. Their keys determine
canonical order and their transitions determine the next state. Continue
by induction. No assumption of a static global alphabet is used.

If the last attempted subset has an internal conflict, do not apply it.
Its candidate catalogue and assignments are still determined by the legal
preceding state. The same slot data uniquely reconstruct that attempted
subset. Thus rejected final attempts are covered without needing a
post-conflict state.
\end{proof}

\begin{lemma}[Exact forest enumeration]
\label{lem:forestcount}
The number of ordered $d$-slot forests with $b$ roots and $j\ge b$ nodes is
\begin{equation}
 F_d(j,b)=\frac{b}{j}\binom{dj}{j-b}.
 \label{eq:forest}
\end{equation}
For $d\ge1$, this is at most $(ed)^j$.
\end{lemma}

\begin{proof}
Let $T(z)$ count nonempty rooted trees by their number of nodes. Choosing
which of the $d$ child slots are filled gives
$T=z(1+T)^d$. A forest of $b$ ordered roots has generating function $T^b$.
The coefficient form of Lagrange inversion yields
\[
 [z^j]T(z)^b=\frac{b}{j}[u^{j-b}](1+u)^{dj},
\]
which is \eqref{eq:forest}. This use of formal power series involves only
finite integer coefficients in each degree.

For the inequality, put $q=j-b$. If $q=0$ the count is one. If $q>0$, use
$\binom{dj}{q}\le(edj/q)^q$ and $b/j\le1$. Writing $t=q/j\in(0,1]$,
\[
 \log\big((edj/q)^q\big)=j\,t\log(ed/t)\le j\log(ed).
\]
Indeed the derivative of $t\log(ed/t)$ is $\log(d/t)\ge0$ on that
interval. This proves the bound, including $d=1$.
\end{proof}

\begin{theorem}[Birth-bounded single-exponential search]
\label{thm:rootbound}
Suppose the initial state has $M$ legal actions and at most $b$ births are
permitted. The layer search through $k$ actions has cost at most a
polynomial factor times
\begin{equation}
 1+\sum_{q=1}^{\min(b,k,M)}\binom Mq
       \sum_{j=q}^{k}\frac{q}{j}\binom{dj}{j-q}.
 \label{eq:exactbudget}
\end{equation}
In particular its cost is
\begin{equation}
 \poly(N,k)\left(1+\sum_{q=1}^{\min(b,k,M)}\binom Mq\right)(ed)^k.
 \label{eq:rootbudget}
\end{equation}
For fixed $d$ and $M=O(N)$, this is $N^{b+O(1)}2^{O(k)}$.
The search can be implemented in polynomial workspace.
\end{theorem}

\begin{proof}
Choose the root subset, charging even subsets that fail the disjointness
test. For every legal root subset use Lemmas~\ref{lem:encoding}
and~\ref{lem:forestcount}. A yielded prefix and an attempted extension are
both bounded by forests of at most $k$ nodes. Candidate enumeration,
subset materialization, compatibility checking, and transitions cost a
polynomial factor per such object under contract (iii). Thus counting only
successful paths does not hide an uncontrolled cost of rejected subsets.
The extra sum over $j$ is absorbed in $\poly(N,k)$.

At most $k$ nonempty layers are active on a depth-first path. Each frame
stores a polynomial-size state, candidate collection, subset iterator, and
history. All integers describing keys, site indices, and counts have
polynomially bounded bit length in the input and $k$. No exponentially
large visited-state cache is required.
\end{proof}

\begin{remark}[An elementary backup bound]
For layer sizes $\ell_1,\ldots,\ell_h$ summing to $j$, a next layer has at
most $d\ell_i$ candidates. Ignoring size compatibility, their subsets are
at most $2^{d\ell_i}$. Multiplying over nonfinal layers gives at most
$2^{dj}$ choices, and the number of compositions of $j$ is at most
$2^{j-1}$. This already proves a $2^{O(dk)}$ bound for fixed roots and
fixed $d$. The forest argument sharpens the base and gives the auditable
coefficient formula; the single-exponential conclusion does not depend on
subtle analytic estimates.
\end{remark}

\section{RIII footprints: 18 darts, nine crossings, 36 successors}
\label{sec:darts}

The diagram representation has four counterclockwise darts at each crossing.
Darts $4c$ and $4c+2$ are the under-strand ports; darts $4c+1$ and $4c+3$
are the over-strand ports. The fixed-point-free involution $\alpha$ pairs
the ends of each edge. Define
\[
 \operatorname{rot}(4c+j)=4c+(j+1\bmod4),\qquad
 \phi(d)=\operatorname{rot}(\alpha(d)).
\]
Faces are cycles of $\phi$. A legal RIII candidate is a triangular face
$T=(d_0,d_1,d_2)$ on three distinct crossings satisfying the over/under
condition: at least one side is over at both ends. Its identity is the
sorted triple of its actual face darts. The crossing triple alone is not
a unique move identifier~\cite{repo-simplify,repo-causal}.

Let $D(T)$ be all twelve darts at its three crossings and let
\begin{equation}
 F(T)=D(T)\cup\alpha(D(T)),\qquad
 V(F)=\{\lfloor d/4\rfloor:d\in F\}.
 \label{eq:footprint}
\end{equation}
This is a conservative read/write footprint. Even when two exterior
strands are directly connected, the formula includes their partner darts.
It avoids an unjustified assumption that the six exterior ports have
six distinct remote neighbors. Write $Q(T)$ for the set of the three core
crossings themselves. The distinction between $Q(T)$ and $V(F(T))$ will
be important: a small core is not a safe replacement for a read/write
footprint.

\begin{lemma}[Local support bounds and closure]
\label{lem:18}
A legal nondegenerate RIII move has $|F(T)|\le18$ and $|V(F(T))|\le9$.
The set $F(T)$ is $\alpha$-closed before the move, is unchanged as a set
by the move, and is closed under the new edge pairing. Every changed
pairing entry lies in $F(T)$. Its crossing support induces a connected
subgraph of the current projection graph.
\end{lemma}

\begin{proof}
The three sides of the triangular face use six distinct darts in $D(T)$,
already paired to one another. Therefore at most the other six darts of
$D(T)$ have partners outside $D(T)$; adding those partners gives at most
$12+6=18$ darts. At the crossing level there are the three central
crossings and at most six exterior endpoint crossings, for a total of
nine. Coincidences only reduce these bounds.

Closure before the move follows immediately from $\alpha^2=1$:
$\alpha(D\cup\alpha(D))=D\cup\alpha(D)$. The move interchanges the
order in which the three strands meet, retaining the twelve central
ports and retaining precisely the same exterior partner darts. It
reconnects pairings among this set, never introducing another exterior
endpoint. Hence its footprint and closure are preserved, and no entry
outside it changes. Equivalently this can be checked directly in the
pairing patch of \code{triangle\_patch}.

The three core crossings form a triangle in the projection graph.
Every additional crossing in $V(F)$ is joined to one of those crossings
by an edge whose dart lies in $D$. Consequently the induced support
subgraph is connected. Loops and repeated edges do not invalidate this
argument.
\end{proof}

\begin{lemma}[Site incidence]
\label{lem:incidence}
At any state, each dart belongs to at most eight candidate RIII
footprints. The total number $M$ of candidates is at most
$\min(N+2,\lfloor4N/3\rfloor)$. If only moves with $V(F)\subseteq S$
are eligible, their number is at most $\lfloor4|S|/3\rfloor$.
\end{lemma}

\begin{proof}
A dart $x$ can lie in $F(T)$ only if $x$ or $\alpha(x)$ lies in $D(T)$.
Thus a candidate must use the crossing of $x$ or the crossing of
$\alpha(x)$. At a fixed crossing, at most four face corners can belong
to triangular faces. Counting each candidate at most once per relevant
corner gives the bound eight.

A connected spherical four-valent projection with $N$ vertices has $N+2$
faces by Euler's formula. Distinct triangular face cycles use disjoint
sets of three darts, so their total number is also at most $4N/3$.
For an eligible candidate supported in $S$, its three face darts all
belong to crossings of $S$. Counting those darts proves the local bound.
Over/under and nondegeneracy tests can only discard candidates.
\end{proof}

Using only these bounds would give $d=18\cdot8=144$. Closure gives a
substantially smaller constant.

\begin{lemma}[Thirty-six next-layer candidates]
\label{lem:36}
After executing a layer of pairwise footprint-disjoint RIII moves, at
most 36 legal RIII candidates intersect the footprint $F$ of any
specified parent move from that layer.
\end{lemma}

\begin{proof}
By Lemma~\ref{lem:18}, $F$ is unchanged and closed under the post-layer
pairing $\alpha$: disjoint sibling moves do not change entries in $F$.
For a candidate with core $D'$ and footprint $F'=D'\cup\alpha(D')$,
\[
 F'\cap F\ne\varnothing\quad\Longleftrightarrow\quad D'\cap F\ne\varnothing.
\]
The reverse implication is immediate. For the forward implication, an
intersection at $\alpha(y)$ with $y\in D'$ implies $y\in F$ by closure.
Such a candidate uses a crossing in $V(F)$, which has size at most nine.
There are at most four triangular-face corners at each such crossing.
Thus at most $9\cdot4=36$ candidates can occur.
\end{proof}

\begin{corollary}[RIII layer budget]
\label{cor:r3layer}
The complete search for RIII traces of at most $k$ actions and at most
$b$ births has cost $N^{b+O(1)}2^{O(k)}$. More precisely,
\eqref{eq:exactbudget} applies with $d=36$ and the initial candidate
count from Lemma~\ref{lem:incidence}.
\end{corollary}

No planar separator theorem is needed for this corollary. The structural
facts used are bounded valence, fixed ports, local commutation, and a
constant-size pairing update. The sharper nine-crossing support will be
used to eliminate the global choice of births.

\section{The backwards cone of one terminal reduction}
\label{sec:cone}

The next step depends on the terminal predicate being local. For a selected
RI test, take all four darts of its crossing and their partners; this
footprint has at most six darts and three crossing vertices. For a selected
RII test, take the eight darts of its two crossings and their partners;
the bigon's two sides already account for four internal darts, leaving at
most four remote endpoints. This footprint has at most twelve darts and
six crossing vertices. Denote it by $F_\rho$. It is closed under the
current pairing and its crossing support is connected. It contains
every pairing entry needed to recognize the selected legal face and its
over/under condition. The test itself is read-only.

\begin{lemma}[Terminal ancestor cone]
\label{lem:cone}
Let $a_1\cdots a_m$ be a legal trace ending with a particular legal local
test $\rho$. Scan the occurrences backwards, starting with active set
$F_\rho$. Retain an occurrence exactly when its footprint intersects the
active set, and then add its footprint to that set. The retained
chronological subsequence is legal from the same input, makes $\rho$
legal, and has the same retained occurrence footprints. The intersection
graph of retained footprints together with $F_\rho$ is connected.
\end{lemma}

\begin{proof}
An omitted occurrence is disjoint from every later retained occurrence
and from the terminal test footprint. It can be moved to the right of
all those retained occurrences by allowed adjacent exchanges. One may
formalize this by repeatedly taking the rightmost omitted occurrence,
commuting it across the retained suffix, and regarding it as occurring
after the test. Earlier omitted occurrences are treated in the same way.
The resulting trace begins with the retained subsequence, followed by the
test; omitted actions occur afterwards and can be dropped. All retained
local data and the test's local data are unchanged by the exchanges.

Each retained footprint intersects the terminal footprint or a footprint
retained later in the backwards scan. Starting with the terminal as a
single connected node, adjoining retained nodes in scan order preserves
connectedness of the intersection graph.
\end{proof}

For a cone containing at most $k$ RIII moves, its crossing union
\begin{equation}
 S=V(F_\rho)\cup\bigcup_{a\text{ retained}}V(F(a))
 \label{eq:region}
\end{equation}
satisfies $|S|\le9k+6$. A naive assertion that this region is connected in
the initial diagram would still require proof: edges change during RIII.
The following invariant is the missing step.

\begin{lemma}[Initial-component confinement]
\label{lem:initialconnect}
Consider a sequence of local rewrites on a fixed graph vertex set. Suppose
each rewrite has a support contained in $S$, induces a connected subgraph
of the current graph, and changes edges only between vertices of its own
support. Then no rewrite creates an edge between different components
of the \emph{initial induced graph} $G_0[S]$. In particular, each rewrite
support lies in one initial component. If a final connected test support
and all rewrite supports have a connected intersection graph and union
$S$, then $G_0[S]$ is connected.
\end{lemma}

\begin{proof}
Partition $S$ into the components of $G_0[S]$. Initially no induced edge
joins different parts. Suppose this remains true before a particular
rewrite. Its connected support must lie in a single part, since a path
in the induced support would otherwise contain an edge joining two
parts. Every changed edge has both endpoints in that support. Therefore
no changed edge joins two parts and the invariant persists. This proves
the first assertion by induction.

The final test support is also connected in the current induced graph,
so it lies in one initial part. Intersecting supports cannot lie in
distinct parts. The connected intersection graph forces all supports
into the same part, and their union is $S$. Thus there is only one part.
\end{proof}

\begin{theorem}[Connected initial support for unlocking]
\label{thm:connectedcone}
If $\mathrm{Unlock}_{\RIII}(D,k)$ is true and the input has no immediate
RI/RII move, it has a witness of at most $k$ RIII moves such that all
move footprints and one terminal reduction footprint have crossing
supports in a set $S$ satisfying
\[
 |S|\le9k+6,\qquad G_D[S]\text{ is connected}.
\]
Here $G_D$ is the initial projection graph, not a graph chosen after
performing the witness.
\end{theorem}

\begin{proof}
Apply Lemma~\ref{lem:cone} to any successful trace and one terminal
reduction. Work with the retained trace from the original input; omitted
moves no longer participate in its state evolution. Lemma~\ref{lem:18}
gives connected crossing supports of size at most nine and changes only
edges with endpoints in those supports. The terminal support has size at
most six and is connected. Intersecting dart supports have intersecting
crossing supports. Their intersection graph is therefore connected.
Lemma~\ref{lem:initialconnect} applies to the union in \eqref{eq:region}.
The size estimate follows by summing the individual support bounds.
\end{proof}

\begin{remark}[Why this is not a geometric disk assumption]
The induced region $G_D[S]$ may have cycles, holes in a planar picture,
or several boundary components. The algorithm keeps the complete diagram
and all external pairings. It does not cut out an unspecified disk or
replace the complement by arbitrary boundary data. The full-support
variant restricts the whole footprint to $S$. The sharper core variant
below uses cores only to restrict which moves may be chosen, while
retaining their entire read/write footprints and the full ambient state.
Neither variant hides an unjustified tangle replacement.
\end{remark}


\subsection{A sharper invariant: connected cores rather than neighborhoods}
\label{sec:cores}

The preceding support theorem is sufficient for single-exponential search,
but it pays to enumerate exterior neighbors that are already determined
by a much smaller set. For any crossing set $C$, define the dart star
\[
 D(C)=\{4c+j:c\in C,\ 0\le j<4\},\qquad
 \mathcal F_\alpha(C)=D(C)\cup\alpha(D(C)).
\]
Its crossing support is exactly the closed neighborhood $N_G[C]$ in the
current projection graph. Here the notation $N_G[C]$ includes $C$ itself.

\begin{lemma}[Core contraction and star invariance]
\label{lem:star}
Let a RIII move have core $Q$, and let $C$ contain $Q$. The move preserves
$\mathcal F_\alpha(C)$ as a set of darts, and it preserves the component
partition of the induced graph $G[C]$. Furthermore, for any two nonempty
core sets $C,Q$ in the same state,
\[
 \mathcal F_\alpha(C)\cap\mathcal F_\alpha(Q)\ne\varnothing
 \quad\Longleftrightarrow\quad
 C\cap Q\ne\varnothing\ \text{or an edge joins }C\text{ to }Q.
\]
\end{lemma}

\begin{proof}
Before and after the move, the three vertices of $Q$ contain the sides of
a triangular face and hence form a connected subgraph. The move retains
the same exterior endpoint darts and reattaches each to a port in $Q$.
Edges with both endpoints outside $Q$ are unchanged. Consequently,
contracting $Q$ to one vertex gives the same graph outside that vertex
before and after the move, with the same incident exterior edges. This
also holds after inducing on any set $C$ containing $Q$. Since $Q$
itself stays connected, the component partition of $G[C]$ is unchanged.

For star invariance, the full footprint $\mathcal F_\alpha(Q)$ is already
contained in both the old and the new $\mathcal F_\alpha(C)$. Its set of
darts is unchanged by Lemma~\ref{lem:18}. For a dart in
$D(C)\setminus\mathcal F_\alpha(Q)$, the pairing is unchanged. A dart
of $D(C)$ inside that footprint has both its old and new partner inside
the same footprint. These observations account for every member of the
old and new stars and prove their equality.

Both stars are pairing-closed. Their intersection is nonempty exactly
when $D(Q)$ meets $\mathcal F_\alpha(C)$. A dart of $D(Q)$ either lies
at a crossing of $C$, or is paired to a dart at a crossing of $C$.
These are precisely the two alternatives in the equivalence.
\end{proof}

\begin{theorem}[Connected initial core of a terminal cone]
\label{thm:core}
Under the hypotheses of Theorem~\ref{thm:connectedcone}, there is a
successful retained trace of at most $k$ RIII moves for which the union
$C$ of the move cores and one terminal reduction core satisfies
\begin{equation}
 |C|\le3k+2,\qquad G_D[C]\text{ is connected}.
 \label{eq:corebound}
\end{equation}
The full dart union of its backwards cone equals
$\mathcal F_{\alpha_0}(C)$ in the initial state. Its crossing support is
therefore $N_{G_D}[C]$; moreover it has at most $7k+6$ vertices.
\end{theorem}

\begin{proof}
Run the backwards selection procedure on a successful chronological
trace, physically considering the successive pre-move states. In
addition to its active dart set $A$, retain the union $C$ of the cores
already selected, beginning with the one- or two-crossing terminal core.
Maintain the following invariant: in the current reverse-scan state,
$A=\mathcal F_\alpha(C)$ and $G[C]$ is connected. It holds at the terminal
state because a monogon core has one crossing and a bigon core has two
adjacent crossings.

Consider a move with core $Q$ and its unchanged footprint
$F=\mathcal F_\alpha(Q)$ in the state immediately after it. If $F$ is
disjoint from $A$, the backwards algorithm omits the move. Reversing it
changes no entry in $A$, and changes no edge incident to a dart of
$D(C)\subseteq A$. Thus both parts of the invariant persist with the
same $A,C$.

If $F$ meets $A$, retain the move. Lemma~\ref{lem:star} says that $Q$
intersects or is adjacent to $C$. The core $Q$ is a triangle also in the
post-move state, so $G[C\cup Q]$ is connected there. Updating the active
set gives
$A\cup F=\mathcal F_\alpha(C\cup Q)$. Now reverse the move. Because
$C\cup Q$ contains the whole move core, star invariance and component
invariance from Lemma~\ref{lem:star} give the desired invariant in the
pre-move state. Induction reaches the initial state. At most three
crossings are added per retained move and at most two are present in
the terminal core, proving \eqref{eq:corebound}. The asserted identity
of the full dart union follows from the same invariant.

For the last estimate, each retained move has a crossing footprint of
size at most nine. At a retained step, two nonempty stars intersect. If
their cores are disjoint, a joining edge places its two distinct endpoint
crossings in the intersection of the crossing supports. If the cores
share a crossing $v$, their star intersection contains the whole star
of $v$. In a connected projection graph with more than one vertex, $v$
has a neighbor distinct from itself, again giving at least two shared
crossings. A nonempty RIII trace necessarily has at least three vertices.
Thus a retained move adds at most seven new crossing vertices to the
active support. Starting with at most six in the terminal support gives
$7k+6$. Zero-length witnesses do not require this argument.
\end{proof}

This proof uses the special exterior-port behavior of RIII, not just
constant footprint size. Arbitrary rewrites inside a closed footprint
could pair two exterior endpoints to one another and disconnect their
old core. Such operations would invalidate the contraction argument;
the theorem must not be transferred to them without checking it.

The proof also explains why unrelated births can be discarded here, but
not for a completely arbitrary global target. A global predicate can
require changes in many disjoint regions without having a constant-size
terminal footprint. Theorem~\ref{thm:rootbound} still applies to that
setting; Theorem~\ref{thm:connectedcone} need not.

\section{An unconditional single-exponential unlocking algorithm}
\label{sec:main}

The simple projection graph obtained by ignoring loops and merging parallel
edges has maximum degree at most four. For nonempty one-component diagrams
it is connected. We need connected \emph{vertex sets}, not induced trees.

\begin{lemma}[Enumerating connected regions]
\label{lem:regions}
In a maximum-degree-four graph with $N$ vertices, the number of nonempty
connected vertex sets of size at most $s$ is $O(N16^s)$. They can be
enumerated without repetition in polynomial workspace and with polynomial
overhead per emitted set.
\end{lemma}

\begin{proof}
For a connected set of size $r$, choose its least vertex as root and a
canonical rooted spanning tree. A depth-first traversal of the tree is a
walk of length $2(r-1)$ in the original graph whose visited vertex set is
exactly the chosen set. A fixed walk identifies its visited set, so mapping
each set to its canonical traversal is injective. There are at most
$N4^{2(r-1)}=N16^{r-1}$ such walks. Summing over $r\le s$ proves the
bound.

For unique enumeration, root a region at its least vertex. Maintain a
selected connected set, an adjacent frontier, and a permanently excluded
set. Choose the least frontier vertex, branch on inclusion or exclusion,
and on inclusion add its eligible neighbors to the frontier. Vertices
below the root are never eligible. Emit the singleton root once and every
new selected set on an inclusion branch. Every connected set has a unique
include/exclude history: its membership dictates each branch while its
connectivity guarantees an available frontier path to each member. The
same selected set is not re-emitted along exclusion-only branches.

A frontier or excluded vertex is adjacent to a selected vertex, so along
a branch there are at most $4s$ such vertices. Every branching node emits
an inclusion child; the full binary search tree has only linear overhead
relative to its inclusion events and leaves. Selected sets, iterator state,
and the explicit depth-first stack therefore use space polynomial in
$N,s$, with polynomial work per output. Stop extending once the selected
size reaches $s$.
\end{proof}

\begin{theorem}[Complete bounded RIII unlocking]
\label{thm:main}
There is a deterministic algorithm for
$\mathrm{Unlock}_{\RIII}(D,k)$ on classical knot diagrams with running time
\begin{equation}
 \poly(N,k)2^{O(k)}
 \label{eq:mainfpt}
\end{equation}
and polynomial workspace. For $k\ge1$, one conservative explicit
upper bound is
\begin{equation}
 \poly(N,k)\,2^{23k+11}.
 \label{eq:explicit}
\end{equation}
It outputs a replayable witness on a positive instance and a complete
negative answer only for this bounded local subproblem.
\end{theorem}

\begin{proof}
First test every RI/RII face. A positive test gives a zero-move witness.
If there is none and $k=0$, return the bounded negative answer. Otherwise
put $t=3k+2$. Enumerate all nonempty connected vertex sets $C$ of size
at most $t$ in the initial projection graph. For each, run the layer
search through $k$ actions, without a birth bound, permitting precisely
those moves whose three-crossing core is contained in $C$. Keep the full
ambient state and continue to use the full dart footprint for legality,
commutation, predecessor incidence, and transitions. Test for a terminal
reduction after each yielded layered prefix.

Every returned path consists of legal RIII moves followed by a legal
test, so the algorithm is sound. For completeness, choose the connected
core $C$ given by Theorem~\ref{thm:core}. All moves of its retained
witness are eligible. Disjoint swaps preserve their cores and full
footprints, so their layer normal form is eligible as well.
Theorem~\ref{thm:foata} guarantees that the search visits a successful
endpoint. If $N\le t$, it suffices to run once on the whole vertex set,
since that set contains every possible witness core.

For complexity, there are $O(N16^t)$ regions by
Lemma~\ref{lem:regions}. A root action eligible for $C$ is a triangular
face with three distinct darts at crossings of $C$. There are at most
$4|C|/3\le4t/3$ such faces. Summing all root subsets costs at most
$2^{4t/3}$. Theorem~\ref{thm:rootbound} with $d=36$ bounds the search
for each region by $\poly(N,k)2^{4t/3}(36e)^k$. Consequently
\begin{equation}
 T(N,k)\le\poly(N,k)\,16^{3k+2}\,
                   2^{(4/3)(3k+2)}(36e)^k.
 \label{eq:product}
\end{equation}
Since $36e<2^7$, the exponential is at most
$2^{12k+8+4k+8/3+7k}\le2^{23k+11}$. Counts of roots, lengths, and
the factor $N$ have been absorbed into the polynomial. The displayed
base is a complexity bound, not a prediction of operational cost.

Process regions one at a time. Their enumeration has polynomial
workspace by Lemma~\ref{lem:regions}; the layer search has polynomial
workspace by Theorem~\ref{thm:rootbound}. This proves the space claim.
\end{proof}

\begin{remark}[Two safe localization modes]
The supplied default region mode uses the connected core bound $3k+2$.
A second, slower theoretical variant enumerates connected full-support
regions of size at most $7k+6$ and restricts whole footprints to those
regions. Both keep the ambient diagram. Using only core crossings as
\emph{read/write footprints} would be unsafe; using them as an additional
eligibility restriction while still checking the full footprint is safe.
\end{remark}

\subsection{Bit complexity and implementation costs}

A dart index uses $O(\log(N+1))$ bits. The state has $4N$ entries, so its
pairing array occupies $O(N\log(N+1))$ bits. Every move changes a constant
number of entries and can be specified by three dart indices. A complete
revalidation of an explicit state, or even a polynomially inefficient
candidate scan, contributes only to the polynomial prefactor in
\eqref{eq:mainfpt}. The bound is consequently not dependent on a
unit-cost assumption for arbitrarily large coordinates or coefficients.

The supplied reference code deliberately does not optimize this prefactor:
it stores immutable tuples, computes faces explicitly, and uses Python
sets and sorted keys. In particular, candidate generation currently may
construct a temporary state to reject a degenerate reconnection. This is
still polynomial work per candidate; removing that copy is a possible
engineering change, not a condition needed for the theorem.

The parameter is the numerical value of $k$. Encoding $k$ in binary does
not turn $2^{O(k)}$ into a polynomial bound in its bit length. The claim
is fixed-parameter tractability with a single-exponential parameter
function, not strong polynomiality in an unbounded numerical budget.

\subsection{Algorithmic skeleton}

\begin{lstlisting}
Unlock(D, k):
    if a legal RI/II face exists in D: return zero-length witness
    if k == 0: return NO_BOUNDED_UNLOCK
    if there is no legal first RIII: return NO_BOUNDED_UNLOCK
    t = 3*k + 2
    regions = connected initial vertex sets of size at most t
    if N <= t: regions = [all vertices]
    for C in regions:
        search legal commuting layers of total size at most k
        keep full ambient diagram; allow a iff its core is in C
        use full F(a) for independence and previous-layer contact
        use every root subset; charge all individual moves
        if a completed prefix exposes RI/II: return its certificate
    return NO_BOUNDED_UNLOCK
# A trial/region/deadline cap returns UNKNOWN, never the last line.
\end{lstlisting}

An empty initial RIII catalogue also permits immediate complete bounded
failure, since no nonempty trace can begin. This shortcut is already
present in the inspected production causal-search discussion~\cite{repo-causal};
it is not a new knot-type criterion.

The phrase ``kernel'' here denotes a bounded-support search component.
We do \emph{not} claim a many-one polynomial kernelization reducing every
$N$-crossing instance to one equivalent $O(k)$-crossing instance. There can
be $N2^{O(k)}$ candidate regions, and external pairing data remain present.

\section{Consequences for the complete recognition pipeline}
\label{sec:pipeline}

The local result becomes relevant to the target complexity only when
combined with an explicit policy and a controlled residual computation.
Existence of a favorable but unselected sequence is not a bound on the
actual implementation.

\begin{definition}[Budgeted deterministic reduction policy]
Fix deterministic choices for RI/RII exhaustion, connected-region order,
layer order, and witness commitment. Start from $D$. Exhaust crossing-
decreasing RI/RII moves. At each stalled nonempty diagram, call the exact
uncapped bounded unlocking algorithm with budget $g$. On success, apply
the returned RIII path and at least one exposed reduction, then repeat
RI/RII exhaustion. On bounded failure, stop preprocessing and send the
residual diagram to an exact fallback. Let $r$ be that residual diagram's
crossing count; put $r=0$ when the empty diagram is reached.
\end{definition}

\begin{lemma}[A suitable exponential fallback exists]
\label{lem:fallback}
A classical knot diagram with $r$ crossings can be recognized exactly in
$2^{O(r)}$ bit time by explicit reduced Khovanov computation over
$\mathbb Q$.
\end{lemma}

\begin{proof}
There are $2^r$ resolutions. Each has at most $r+1$ circles, so the
number of basis generators in the unreduced complex is at most
$2^{2r+1}$; reducing does not increase it. The usual merge and split maps
have entries $0,1,-1$. Explicit construction is bounded by a polynomial
in $r$ and this exponential dimension. Fraction-free elimination computes
the ranks of all differential matrices. Hadamard's determinant bound
shows that their minors have bit length polynomial in the dimension
and in $r$, hence $2^{O(r)}$. Polynomially many arithmetic operations
in that dimension with these bit lengths still cost $2^{O(r)}$.
The total rational homology rank is then exact. Reduced rational
Khovanov homology has rank one exactly for the unknot by the
unknot-detection theorem~\cite{km}.
\end{proof}

This lemma supplies a complexity-compatible complete fallback; it does
not say that a new homology implementation is included in this package.
The repository already has complete exact backends. The new package
contains only the local search, its verification machinery, and a proposed
bridge to its mutable dart representation.

\begin{theorem}[Run-dependent gap/core bound]
\label{thm:gapcore}
For the deterministic reduction policy above, using an exact fallback
with cost $2^{O(r)}$, recognition has cost
\begin{equation}
 \poly(N,g)2^{O(g)}+2^{O(r)}.
 \label{eq:gapcore}
\end{equation}
In particular, a family on which the \emph{actual run} has
$g,r=O(\log^2 N)$ admits recognition in $N^{O(\log N)}$ time.
If the policy reaches the empty diagram, it produces an unknot witness
of $O(N(g+1))$ Reidemeister moves.
\end{theorem}

\begin{proof}
Each successful episode includes a crossing-decreasing move. There are
at most $N$ such episodes, and the number of crossings during every RIII
search is at most $N$. RI/RII exhaustion is polynomial in the explicit
input size; even rebuilding after each reduction suffices for that claim.
Theorem~\ref{thm:main} gives at most
$N\poly(N,g)2^{O(g)}$, with the factor $N$ absorbed into the polynomial.
There is at most one unsuccessful final bounded call, and its cost has
the same bound. Add the residual fallback cost.

If $g,r\le C\log^2N$, the exponential factors are
$2^{O(\log^2N)}=N^{O(\log N)}$; the polynomial prefactor does not
change this form. There are at most $N$ successful episodes, each with
at most $g$ RIII moves and a total of at most $N$ crossing-decreasing
moves over the entire run. This proves the certificate-length bound.
\end{proof}

With fixed original dart labels retained across deletions, the positive
certificate uses $O(N(g+1)\log(N+1))$ bits. A verifier replays the moves
without searching. Full revalidation after every move is polynomial in
$N$ and certificate length; incremental verification can be faster but
is not needed for this assertion.

\subsection{Where the universal theorem is still missing}

Nothing here proves that every unknot admits a successful RIII-only
unlocking episode with $g=O(\log^2N)$ at each stall of this policy.
Nothing proves that the residual $r$ is polylogarithmic when the bounded
search fails. The method never authorizes a knottedness verdict merely
from local failure, even after complete bounded exhaustion.

These are substantive geometric and policy questions, not a missing
constant in the combinatorial analysis. Allowing crossing-increasing moves
would require new support bounds, dynamic creation and deletion of ports,
and corresponding completeness arguments. A universal hierarchy method
would require its own representations and costs; it cannot be inferred
from the local causal cone.

The improvement nevertheless changes a useful scale. The bounded local
query is polynomial for $k=O(\log N)$ and fits the specified
$N^{O(\log N)}$ budget for $k=O(\log^2N)$, without any independent-birth
promise. The old $k^k$ factor at the latter scale contributes an additional
$\log\log N$ in the exponent. Root removal additionally avoids a large
$N^b$ term when many births are possible.

\section{Executable design, certificates, and integration safety}
\label{sec:implementation}

\subsection{Package components}

\begin{center}
\begin{tabular}{>{\raggedright\arraybackslash}p{0.32\textwidth}>{\raggedright\arraybackslash}p{0.60\textwidth}}
\toprule
Module & Responsibility\\
\midrule
\code{layered\_search.py} & Generic pure-state layer enumeration, bounded roots,
normal-form and backwards-cone utilities, explicit resource exceptions.\\
\code{dart\_kernel.py} & Classical diagram validation, local RIII pairing
patches, RI/RII tests, and reference reduction operations.\\
\code{connected\_kernel.py} & Unique connected-region enumeration and the
complete bounded unlocking algorithm.\\
\code{verify\_certificate.py} & Search-independent certificate replay and
validation of layers and terminal faces.\\
\code{production\_adapter.py} & Snapshot/commit bridge retaining original
production dart identities and rolling back on failure.\\
\code{unlock.py} & Command-line bounded query with explicit
\code{FOUND\_UNLOCK}, \code{NO\_BOUNDED\_UNLOCK}, and \code{UNKNOWN} statuses.\\
\code{audit.py}, \code{benchmark.py} & Deterministic mathematical audits and
paired kernel timings, with raw JSON output.\\
\bottomrule
\end{tabular}
\end{center}

All Python code uses the standard library. No Lean formalization or other
proof-assistant artifact is claimed. The certificate verifier is independent
of the \emph{search control flow}, but shares the dart reconnection kernel.
Thus it is not an independent second implementation of knot topology.
The comparisons against chronological enumeration independently exercise
the search logic, not the provenance of the common Reidemeister formulas.

\subsection{One global copy for a disjoint layer}

For a pure-state implementation, copying the $4N$-entry pairing array after
each move costs $O(N\ell)$ for a layer of width $\ell$. Since footprints
are disjoint, compute each constant-size patch from the common pre-layer
state, allocate one copy, and apply the patches to it. The cost becomes
$O(N+\ell)$ pairing-entry operations, after candidate validation. The
implemented \code{apply\_layer} uses precisely this method. It does not
claim a corresponding speedup over the production mutable journal, which
already avoids global copies in a different way.

If patch creation or validation fails, the original immutable state is
untouched. An exception while constructing a tentative result cannot leak
a partially modified input. This is useful for resource-safe research
experiments even when the copying overhead is not competitive.

\subsection{Certificate contract}

A certificate records the input PD code, a maximum move count, a sequence
of nonempty layers of three-dart face keys, and a specific terminal RI/RII
face. Its format identifier is
\code{proveit-r3-layer-certificate-v1}. The verifier checks classical
one-component input validity, literal integer types, legal face identity,
canonical key ordering, pairwise footprint disjointness in a layer,
intersection with the preceding layer, legal transitions, move count, and
the selected terminal reduction. It revalidates the diagram after every
move.

The face key matters. A set of three crossing indices can identify two
different triangular faces, as the repository's explicit audit
shows~\cite{repo-causal}. The new format identifies darts in the current
fixed-label episode, not just crossings. The bridge maps those compact
episode labels back to the original live production dart labels after
previous crossing deletions.

Certificates are positive witnesses only. A completed search can provide
a trustworthy local negative \emph{computation}, but we do not include a
succinct independently checkable certificate of bounded nonreachability.
A timeout or exhausted subset/region allowance cannot be serialized as a
negative certificate.

\subsection{Safe opt-in adapter}

The adapter snapshots the production object's live crossings and pairing
entries, validates the resulting compact diagram, and records the mapping
to original labels. Search occurs on the isolated snapshot. Before any
commit, the candidate certificate is replayed, its final pairing is checked,
and the live production object is compared with the original snapshot to
reject stale data.

The bridge reconstructs actual cyclic faces from verified keys; it does
not trust an action's optional payload. It journals the original live
pairings and the original path length before writing. Any exception,
including a cooperative cancellation during mutation, restores the
pairing and truncates the path to that original length. It does not
perform the exposed crossing deletion, change production trial counters,
change defaults, or invent a verdict.

This is an integration proposal with unit-tested invariants, not an
already integrated production patch. The full repository regression
suite was not run with this bridge. Compatibility with evolving internal
APIs, budget accounting, restart handling, and the existing independent
trace replayer must be checked in the real checkout before enabling it.

\section{Validation and measured performance}
\label{sec:experiments}

All results in this section are completed runs of the delivered code.
The environment was CPython 3.13.5 on Linux x86-64, with platform string
recorded in the raw benchmark JSON. No performance result here is a
measurement of the complete \code{fastunknot} recognition pipeline.

\subsection{Tests and exhaustive checks}

The final local unit suite has 27 passing tests. It includes invalid
parameters, caps, cancellation, the inverse-padding counterexample,
canonical fork/join layers, connected-set enumeration, actual RIII face
ambiguity, disjoint commutation, batch-versus-sequential transitions,
certificate tampering, deleted-label mapping, stale snapshots, rollback,
and rejection of a forged final state. The adapter also has a regression
test ensuring that an untrusted action payload cannot be committed as
an unchecked move identity.

The independent Boolean audit uses all 27 guarded one-bit toggles on
three sites: choose the toggled bit and independently choose whether each
other bit is untested, required zero, or required one. Each rule is an
involution because its guard does not read its flipped bit. All rule sets
of sizes one through three give
$\binom{27}{1}+\binom{27}{2}+\binom{27}{3}=3303$ systems. For all eight
initial states and birth bounds one through three, compare chronological
enumeration through depth four with the layered endpoint set. That gives
79,272 comparisons, covering 1,313,568 reference trace visits, with zero
discrepancies. Equality of endpoint sets is tested, not merely equality
of a yes/no outcome on a few selected goals.

The forest formula is checked independently by truncated series
convolution of $T=z(1+T)^d$ rather than by reusing its binomial expression.
There are 651 coefficient comparisons over the recorded ranges, with
zero discrepancies. These finite comparisons support the implementation;
Lemma~\ref{lem:forestcount}, not the tests, proves the formula for all
parameters.

\subsection{Actual-diagram audit}

Using seed 8102601, the audit generates 1,200 braid words on three through
six strands, with lengths eight through 32. It retains the 248 words
whose closures validate as knots, then performs deterministic RI/RII
exhaustion. On eligible resulting states it checks individual RIII
supports and moves, disjoint pairs, and bounded endpoint sets. It also
runs random legal RIII traces of length at most eight, extracts the
backwards cone of available terminal reductions, replays the retained
subsequence and its layer form, and checks both support and core connectivity in the initial
projection graph. The same 1,000 cases check the exact identity of the
cone's full dart union with the initial dart star of its retained core
union, and the corresponding closed-neighborhood identity on crossings.

\begin{center}
\begin{tabular}{lr}
\toprule
Audit item & Completed count\\
\midrule
Validated knot closures & 248\\
Bounded actual-diagram endpoint comparisons & 180\\
RIII action/support/inverse checks & 821\\
Disjoint actual-diagram commutation checks & 543\\
Terminal-cone extraction and initial-connectivity checks & 1,000\\
Core connectivity and initial-star identities (same cones) & 1,000\\
Nonempty cones among those checks & 1,000\\
Retained positive certificate files & 10\\
Discrepancies in these checks & 0\\
\bottomrule
\end{tabular}
\end{center}

The largest observed footprint has 18 darts and nine crossing vertices,
attaining both proved local upper bounds. The largest observed site
incidence is five and the largest observed next-parent candidate count
is nine; these smaller observations do not replace the uniform proved
bounds eight and 36. The ten certificates in \code{artifacts/} are
nonempty positive witnesses. They demonstrate replay, not difficult
unknot recognition or a lower bound on minimal unlocking depth.

\subsection{Paired kernel timings, including negative results}

We compare a pure-state implementation of the old birth-front normal
form with the new commuting-layer enumerator. Both use the same transition
kernel; the baseline omits the production code's optional inverse pruning.
Every trace prefix through the specified depth is exhausted. Goal tests
do not terminate these timing runs early. This measures interleaving
work, not the time to find a successful reduction, and not full
recognition. An identical baseline control is timed in each round.
There is one untimed warmup of the distinct algorithms, followed by five
paired rounds with randomized arm order, seed 8102602.

The knot cases are the closures of $(\sigma_1\sigma_2)^q$ for
$q\in\{4,7,10\}$. Their purpose is to provide genuine diagram operations
with controlled independent regions, not to represent a hard-knot corpus.
The six-bit cases are explicitly abstract systems. The paired ratio in
Table~\ref{tab:bench} is the median of baseline time divided by layer time
\emph{within each round}; it is not in general the quotient of the two
reported marginal medians.

\begin{table}[htbp]
\centering\small
\begin{tabular}{lrrrrrr}
\toprule
Case & $k$ & Old traces & Layers & Old ms & New ms & Paired ratio\\
\midrule
Six bits & 1 & 6 & 6 & 0.010 & 0.018 & 0.549\\
Six bits & 6 & 1,653 & 923 & 4.526 & 4.614 & 0.952\\
Six bits & 8 & 15,327 & 3,002 & 35.911 & 15.747 & 2.261\\
$(\sigma_1\sigma_2)^4$ & 1 & 8 & 8 & 0.102 & 0.109 & 0.925\\
$(\sigma_1\sigma_2)^4$ & 4 & 664 & 664 & 8.918 & 9.963 & 0.896\\
$(\sigma_1\sigma_2)^7$ & 4 & 6,069 & 4,928 & 85.497 & 86.583 & 0.980\\
$(\sigma_1\sigma_2)^{10}$ & 4 & 18,095 & 14,305 & 291.492 & 307.004 & 0.952\\
\bottomrule
\end{tabular}
\caption{Exhaustive kernel timings. ``Layers'' counts yielded layered
trace prefixes, not individual layer objects. Ratios above one favor the
new code. These are not end-to-end recognition speedups.}
\label{tab:bench}
\end{table}

For the 14- and 20-crossing diagram cases, the new normal form visits about
18.8\% and 20.9\% fewer trace prefixes, respectively. However, layer
construction and subset testing erase that advantage in the measured
wall-clock times. The actual-diagram paired ratios range from 0.896 to
0.980: there is no demonstrated knot-kernel speedup in these runs. The
abstract six-bit depth-eight case does improve, but promoting its
$2.261$ ratio to a recognition claim would be misleading.

The identical-control paired ratios range from approximately 0.963 to
1.024 across these cases. This is a small observational noise check, not
a confidence interval or a statistical performance guarantee. Very short
measurements, especially the depth-one Boolean row, are intrinsically
sensitive to scheduling and timer variation. All per-round times and arm
orders are retained in \code{data/benchmark.json}; counts are deterministic.

The study therefore supports a theoretical improvement and a reduction
of some duplicate work, \emph{not} a default implementation switch. The
connected-region algorithm itself has not been subjected to a large
end-to-end performance comparison. Its conservative constant and region
multiplicity may dominate at practical parameter values.

\section{Limits, proof obligations, and failure modes}
\label{sec:limits}

\paragraph{A local terminal is essential.}
The backwards-cone proof starts with one fixed-size terminal footprint.
It cannot be used without modification for a global invariant comparison,
full homology calculation, or a target depending on many separated
regions. In those cases a birth or terminal-support parameter may remain.

\paragraph{The full footprint is essential.}
The three flipped crossings alone are not a safe read/write support.
Partner darts can be read or changed. The core-region algorithm uses
core inclusion only to restrict its action catalogue; every independence,
legality, and update check still uses $D\cup\alpha(D)$, including all
remote endpoint darts. Treating the core as the footprint would destroy
the proof.

\paragraph{Graph evolution must be controlled.}
Being connected in the graph immediately before each move does not, by
itself, imply that a union is connected initially. The confinement lemma
also requires that every changed edge have both endpoints inside the
move's connected support. It is proved for the retained trace, after
irrelevant actions are removed; ignoring that order of reasoning leaves
a gap.

\paragraph{A length bound is not a layer-count bound.}
Charging every nonempty parallel layer one unit could hide arbitrarily
many actions. All statements and code charge the sum of layer sizes.
The certificate verifier checks that sum independently.

\paragraph{Finite caps change the specification.}
The complete theorems assume the search exhausts its domain when no
witness exists. In an operational recognizer, resource caps should be
expected. They produce \code{UNKNOWN} and preserve the existing complete
fallback. The package's \code{SearchExhausted} exception is not a negative
recognition result.

\paragraph{Local choices need not compose optimally.}
Selecting one available reduction can change later unlocking depths.
The gap/core theorem explicitly concerns the actual deterministic policy.
It cannot substitute the existence of a different globally favorable
reduction sequence. The distinction is also consistent with the subtleties
of greedy moves in global untangling problems~\cite{lrt}.

\paragraph{Testing is not proof or complete integration.}
The tests check concrete implementations and many instances. They neither
prove all geometrical claims nor demonstrate compatibility with every
current production pathway. The mathematical claims are justified by the
preceding proofs. The package does not report a rerun of the repository's
full test suite and does not report a new complete recognizer.

\paragraph{Priority and external review.}
The layer machinery builds on classical partial commutation. The connected
initial-support theorem and its combination with the dart bounds are
presented as this contribution's main advances relative to the inspected
repository. The literature review is not a proof that no equivalent
parameterized local-rewrite theorem has appeared elsewhere. The arguments
should receive independent mathematical review before being described as
an established breakthrough.

\section{Further research questions and concrete work packages}
\label{sec:future}

\subsection{Reduce the single-exponential base without weakening completeness}

\textbf{Question 1: Can connected heaps be enumerated directly?}
The factor $16^{3k+2}$ comes from choosing a connected core region before
searching its layers. Different regions can contain the same causal
trace. A direct exploration of connected terminal-rooted dependency
structures may avoid this multiplicity. The difficulty is that the
terminal face is not necessarily a face in the initial diagram. A useful
result would be an injective code for terminal-rooted dynamic heaps with
$N C^k$ codes and a practically smaller $C$, including an explicit decoder.

\textbf{Question 2: Is the next-parent bound 36 sharp?}
The estimate counts four corners at each of nine support crossings and
ignores triangle compatibility, over/under parity, and the fact that six
of the vertices are exterior neighbors. Determine the maximum number of
legal footprints meeting one post-flip footprint in a realizable
one-component spherical map. An exhaustive local rotation-system catalogue
could conjecture a better value, but a uniform proof must handle all
external identifications rather than only generic six-ended tangles.

\textbf{Question 3: Can region enumeration exploit triangle coverage?}
The core union has at most $3k+2$ crossings, but it is not an arbitrary
connected set: it is covered by the cores of retained triangles and one
terminal core. Encode that additional structure directly, while allowing
the triangles to become faces only after earlier operations. Determine
whether this yields fewer candidates than the generic $16^{3k+2}$ region
bound. Bounds based only on a static catalogue of initial triangular
faces would miss dynamically created faces.

\subsection{Turn reduced interleavings into measured speed}

\textbf{Question 4: Which subset generator is best for the current dart
engine?} The reference generator enumerates combinations and then checks
compatibility. A prefix-pruned independent-set generator in the conflict
graph could reject overlaps earlier. A production benchmark must compare
against the existing journal-based search, preserve inverse semantics,
and include all candidate-generation overhead. Node counts alone are not
a sufficient acceptance criterion.

\textbf{Question 5: Can certified state sharing reuse work between
regions?} Two regions often induce identical eligible actions along a
prefix. Cache keys might include the live pairing on an explicitly
closed support, the previous layer, remaining depth, and admissibility
information. Prove a sufficient key before sharing results. Physical
endpoint equality by itself is not enough for every restricted search.
A cache should also have a polynomial-space mode or a separately stated
space bound.

\textbf{Question 6: Can layer batching be implemented with the mutable
production journal?} The current reference code makes one global copy
per layer. The existing production engine already journals local writes.
A joint layer transaction could register all disjoint patches once,
perform a localized terminal scan, and roll back atomically. Required
tests include injected exceptions at every mutation point, deleted-dart
identities, and replay using the repository's existing verifier.

\subsection{Bridge the local theorem to stronger recognition classes}

\textbf{Question 7: What families have polylogarithmic policy gap and
residual core?} Measure the actual quantities entering
Theorem~\ref{thm:gapcore}, not only successful traces found by a different
policy. Natural first studies include diagrams with controlled braid
run structure, small connected-sum cores, and the project's stored hard
unknot survivors. A theorem should specify the chosen normalizations and
reduction order, and should bound the total preprocessing that obtains
its promised representation.

\textbf{Question 8: Can limited crossing increases be incorporated?}
Introduce an additional parameter for created crossings or for a bounded
height above the initial crossing count. Track provenance of new ports
and local terminal ancestors. Does an initial-support argument survive
when graph vertices are created, perhaps by attaching them to the old
ports that created them? A successful result must bound both the number
of created sites and the encoding of their attachment histories.

\textbf{Question 9: Can a policy-independent residual theorem be proved?}
Instead of requiring every selected episode to be favorable, ask whether
all maximal sequences of reductions with a specified local completeness
budget leave a controlled core. Alternatively, identify a efficiently
verifiable potential that makes one policy provably safe. Greedy
crossing count alone does not provide such a theorem.

\textbf{Question 10: Is there a useful terminal-support hierarchy?}
The cone lemma permits a terminal support larger than a single RI/RII
move if its size is explicitly parameterized. One could target a
certified braid-kernel shortening, a small disk simplification, or a
structured multi-move certificate. The intended target must be recognized
locally with controlled support; importing a global invariant oracle
without its cost defeats the bound.

\subsection{Relate local search to hierarchy and complexity theory}

\textbf{Question 11: Can dependency layers organize hierarchy
operations?} Lackenby's hierarchy approach is not a sequence of the
fixed-site RIII moves used here. Nevertheless, one may ask whether an
encoded hierarchy admits certified independent operations with bounded
interaction neighborhoods. Before transferring the forest count, prove
an appropriate representation-size bound, conservative read/write
supports, dynamic commutation, and the analogue of $d$. Merely drawing
an analogy with parallelism would not supply an algorithm.

\textbf{Question 12: What parameterized lower bounds match this local
problem?} Study the precise bounded unlocking language, with no crossing
increases and only one local terminal reduction. Determine whether
$2^{o(k)}\poly(N)$ is plausible under standard hypotheses. Reductions
must preserve classical one-component planarity and the exact allowed
moves. Existing hardness for whole-diagram defect does not automatically
supply this lower bound.

\textbf{Question 13: Can the proofs be formalized modularly?}
A useful formalization order is: finite-site exchange contracts; occurrence
heights and linear extensions; bounded-slot forest injection; finite graph
component confinement; dart footprint lemmas; then the full algorithmic
count. This separates generic combinatorics from Reidemeister topology.
The certificate format and finite audit catalogue provide concrete test
objects, not replacements for these formal proofs.

\subsection{Immediate integration decision}

Retain the current production default. Add the new search as an explicitly
experimental bounded-unlocking mode after replay and budget integration
are verified in a complete checkout. Measure paired end-to-end recognition
on the project's actual hard corpus, with unchanged complete fallbacks,
identical global budgets, randomized arm order, and negative results
retained. Promote a default only on that evidence; preserve the proved
complete mode separately from heuristic early dispatch.

\section{Conclusion}

Two independently useful effects emerge. Commuting layers remove the
$k^k$ interleaving overhead from the birth-bounded local search. A terminal
ancestor cone, together with initial-component and core-contraction
invariants, then removes the birth-choice factor for the specific RI/RII unlocking problem. The
result is a complete deterministic single-exponential-parameter search,
with explicit local constants and polynomial workspace.

The new bottleneck is therefore not enumeration of all chronological
interleavings of a short unlocking episode. It is the geometric size of
an episode needed by the selected simplification policy, the size of the
residual core when that policy stalls, and practical constant factors.
The gap/core theorem states exactly what would suffice to reach
$N^{O(\log N)}$ for the complete pipeline. The included code and audits
establish a reproducible starting point for that next step without
misreporting a general recognition theorem or an unmeasured speedup.

\appendix
\section{Reproduction and integration instructions}

From the package root, a Unix-like shell can run:
\begin{lstlisting}
export PYTHONPATH=src
python -m unittest discover -s tests -v
python src/audit.py
python src/benchmark.py --rounds 5
python src/verify_certificate.py artifacts/witness_00.json
python src/unlock.py artifacts/witness_00.json --depth 4 \
    --max-trials 100000 --certificate /tmp/replayed-unlock.json
sh build.sh
\end{lstlisting}
The input accepted by \code{unlock.py} is either a PD object or an Artin
braid object. A certificate file includes a PD field and is therefore also
a valid query input. The example above searches again; it does not assume
the certificate's listed path is minimal. Omit caps only when complete
bounded exhaustion is intended. A zero cap means zero allowance, not an
unlimited search. Running audits or benchmarks regenerates their JSON
outputs, including local timings; the shipped outputs record the run
reported in this article.

\code{README.md} provides a package map and status summary.
\code{INTEGRATION.md} gives the production adapter contract and the
recommended opt-in integration sequence. \code{SOURCES.json} records
source paths, revisions, blob hashes, and adaptation scope. The archive
contains no third-party preprint PDF or font file, and no complete
production checkout. \code{MANIFEST.sha256} covers the delivered files.
The article is self-contained LaTeX with an inline bibliography, so no
external bibliography download is needed to rebuild it.

\section{Proof-dependency and evidence ledger}

\begin{longtable}{>{\raggedright\arraybackslash}p{0.28\textwidth}>{\raggedright\arraybackslash}p{0.37\textwidth}>{\raggedright\arraybackslash}p{0.25\textwidth}}
\toprule
Statement & Mathematical dependency & Computational evidence\\
\midrule
\endhead
Dynamic normal form & Finite occurrence poset and stable disjoint exchange
contracts & Boolean and diagram endpoint comparisons\\
Forest budget & Deterministic parent/slot assignment; formal coefficient
extraction & 651 convolution comparisons\\
18-dart, nine-crossing bound & Six internal triangular-side darts and
unchanged exterior endpoints & 821 action/support checks\\
36 successor bound & Post-layer pairing closure and four corners per crossing
& Observed maximum nine; not used as proof\\
Terminal cone & Backwards support test and exchanges of omitted actions
& 1,000 retained-trace replays\\
Initial connectedness & Component confinement for full supports; core
contraction and star invariance for connected $3k+2$ cores & 1,000 core
connectivity and initial-star equality checks\\
Single-exponential unlocking & Region count, local root catalogue, forest
budget & Executable exact region/layer search\\
Gap/core recognition bound & At most $N$ successful episodes and an exact
$2^{O(r)}$ fallback & No uniform gap/core measurements or proof claimed\\
Production adapter & Snapshot validation, verified keys, stale-state check,
journaled commit & Local unit tests only; full production suite not run\\
\bottomrule
\end{longtable}

\begin{thebibliography}{99}
\bibitem{repo-causal}
V. Reshetnikov's ProveIt repository,
\emph{Causal RIII search: review and unambiguous certificates},
\code{Topology/UnknotRecognition/synthesis/causal\_r3.tex}.
Inspected 8 October 2026; blob
\nolinkurl{b5b7b3f01db47d358df2a7b9889865c635a42044}.
\repo.

\bibitem{repo-code}
ProveIt contributors, \code{fastunknot/causal\_r3.py},
\code{Topology/UnknotRecognition/fast/}, blob
\nolinkurl{19599a59b461f74c0151acb40c3658a6a61656b7}.
Inspected 8 October 2026. \repo.

\bibitem{repo-simplify}
ProveIt contributors, \code{fastunknot/simplify.py}, inspected lines 1--310
at revision \nolinkurl{34f2689e96e04a13a32d133173b5e41a34320086}; blob
\nolinkurl{703975603890b0481a31d7368502f48fdda22375}.
\repo.

\bibitem{repo-diagram}
ProveIt contributors, \code{fastunknot/diagram.py}, inspected lines 1--220
at revision \nolinkurl{34f2689e96e04a13a32d133173b5e41a34320086}; blob
\nolinkurl{8717f530cc2847492858da46ac4ac40a2288b7ee}.
\repo.

\bibitem{viennot}
G. X. Viennot, \emph{Heaps of pieces, I: Basic definitions and combinatorial
lemmas}, in \emph{Combinatoire \'enum\'erative}, Lecture Notes in Mathematics
1234, Springer, 1986.
\url{https://doi.org/10.1007/BFb0072524}.

\bibitem{krattenthaler}
C. Krattenthaler, \emph{The theory of heaps and the Cartier--Foata monoid},
author manuscript.
\url{https://www.mat.univie.ac.at/~kratt/artikel/heaps.pdf}.

\bibitem{lrt}
C. Legrand-Duchesne, A. Rai, and M. Tancer,
\emph{Parameterized complexity of untangling knots},
arXiv:2111.05001v1, 2021.
\url{https://arxiv.org/abs/2111.05001}.

\bibitem{lackenby}
M. Lackenby, \emph{Incompressible surfaces, hierarchies and unknot
recognition}, arXiv:2607.23350v1, 25 July 2026, especially Section 9.
\url{https://arxiv.org/abs/2607.23350}.

\bibitem{km}
P. B. Kronheimer and T. S. Mrowka, \emph{Khovanov homology is an
unknot-detector}, Publications Math\'ematiques de l'IH\'ES 113 (2011),
97--208. \url{https://doi.org/10.1007/s10240-010-0030-y}.

\bibitem{bn}
D. Bar-Natan, \emph{Fast Khovanov homology computations},
Journal of Knot Theory and Its Ramifications 16 (2007), 243--255;
arXiv:math/0606318.
\url{https://www.math.utoronto.ca/~drorbn/papers/FastKh/}.
\end{thebibliography}
\end{document}
